% !TeX root = ../Book1.tex
\begin{figure}[htbp]
	\centering
	\begin{tikzpicture}[font=\small,x=0.95cm,y=1cm]
		\node[left] at (-0.3,1) {\(V_4\)};
		\foreach \x/\m in {0/4,2/2,4/0,6/-2,8/-4}{
			\fill[GeometryBlue] (\x,1) circle (2pt);
			\node[above] at (\x,1.13) {\(\m\)};
		}
		\foreach \x in {0,2,4,6}{
			\draw[->,GeometryBlue,thick] (\x+0.15,1)--(\x+1.85,1) node[midway,below] {\(F\)};
		}
		\node[right] at (8.35,1) {下降到 \(\operatorname{SO}(3)\)};
		\node[left] at (-0.3,-0.7) {\(V_3\)};
		\foreach \x/\m in {1/3,3/1,5/-1,7/-3}{
			\fill[GeometryBlue] (\x,-0.7) circle (2pt);
			\node[above] at (\x,-0.57) {\(\m\)};
		}
		\foreach \x in {1,3,5}{
			\draw[->,GeometryBlue,thick] (\x+0.15,-0.7)--(\x+1.85,-0.7) node[midway,below] {\(F\)};
		}
		\node[right] at (8.35,-0.7) {\(-I\) 作用为 \(-1\)};
	\end{tikzpicture}
	\caption[特殊酉二阶群的权链]{\(V_4\) 与 \(V_3\) 的每个权空间都是一维。降算子 \(F\) 每次把权减二，升算子 \(E\) 沿相反方向作用；图中箭头对应非零作用，\(F\) 在最低权处为零，\(E\) 在最高权处为零。偶数权链使中心元 \(-I\) 作用平凡，奇数权链则不满足下降条件。}
	\label{g1:fig:su2-weight-chains}
\end{figure}
